\documentclass[10pt,a4paper]{amsart}
\usepackage{amsmath,amssymb,amsthm,graphicx}
% Retained preamble for the Stacks Project chapter divisors.tex (Divisors).
% The chapter file pinned-original-source.tex carries only \input{preamble} on its first line, so no
% preamble text lives inside the chapter itself.  The Stacks Project uses one single shared preamble for
% every chapter; the snapshot stages it verbatim in sources/stacks-chow/preamble.tex.  The block below is
% copied mechanically, line for line, from that file's lines 14-18 (the Xy-pic setup) and lines 154-217
% (the native theorem declarations together with the mathematical macros).  Nothing was added, removed or
% redefined.  The wrapper class and the amsmath/amssymb/amsthm/graphicx base packages are supplied by the
% assembler, so the class/\externaldocument part of the shared preamble is not repeated here.
% For commutative diagrams we use Xy-pic
\usepackage[all]{xy}

% We use 2cell for 2-commutative diagrams.
\xyoption{2cell}
% Theorem environments.
%
\theoremstyle{plain}
\newtheorem{theorem}[subsection]{Theorem}
\newtheorem{proposition}[subsection]{Proposition}
\newtheorem{lemma}[subsection]{Lemma}

\theoremstyle{definition}
\newtheorem{definition}[subsection]{Definition}
\newtheorem{example}[subsection]{Example}
\newtheorem{exercise}[subsection]{Exercise}
\newtheorem{situation}[subsection]{Situation}

\theoremstyle{remark}
\newtheorem{remark}[subsection]{Remark}
\newtheorem{remarks}[subsection]{Remarks}

\numberwithin{equation}{subsection}

% Macros
%
\def\lim{\mathop{\mathrm{lim}}\nolimits}
\def\colim{\mathop{\mathrm{colim}}\nolimits}
\def\Spec{\mathop{\mathrm{Spec}}}
\def\Hom{\mathop{\mathrm{Hom}}\nolimits}
\def\Ext{\mathop{\mathrm{Ext}}\nolimits}
\def\SheafHom{\mathop{\mathcal{H}\!\mathit{om}}\nolimits}
\def\SheafExt{\mathop{\mathcal{E}\!\mathit{xt}}\nolimits}
\def\Sch{\mathit{Sch}}
\def\Mor{\mathop{\mathrm{Mor}}\nolimits}
\def\Ob{\mathop{\mathrm{Ob}}\nolimits}
\def\Sh{\mathop{\mathit{Sh}}\nolimits}
\def\NL{\mathop{N\!L}\nolimits}
\def\CH{\mathop{\mathrm{CH}}\nolimits}
\def\proetale{{pro\text{-}\acute{e}tale}}
\def\etale{{\acute{e}tale}}
\def\QCoh{\mathit{QCoh}}
\def\Ker{\mathop{\mathrm{Ker}}}
\def\Im{\mathop{\mathrm{Im}}}
\def\Coker{\mathop{\mathrm{Coker}}}
\def\Coim{\mathop{\mathrm{Coim}}}

% Boxtimes
%
\DeclareMathSymbol{\boxtimes}{\mathbin}{AMSa}{"02}

%
% Macros for moduli stacks/spaces
%
\def\QCohstack{\mathcal{QC}\!\mathit{oh}}
\def\Cohstack{\mathcal{C}\!\mathit{oh}}
\def\Spacesstack{\mathcal{S}\!\mathit{paces}}
\def\Quotfunctor{\mathrm{Quot}}
\def\Hilbfunctor{\mathrm{Hilb}}
\def\Curvesstack{\mathcal{C}\!\mathit{urves}}
\def\Polarizedstack{\mathcal{P}\!\mathit{olarized}}
\def\Complexesstack{\mathcal{C}\!\mathit{omplexes}}
% \Pic is the operator that assigns to X its picard group, usage \Pic(X)
% \Picardstack_{X/B} denotes the Picard stack of X over B
% \Picardfunctor_{X/B} denotes the Picard functor of X over B
\def\Pic{\mathop{\mathrm{Pic}}\nolimits}
\def\Picardstack{\mathcal{P}\!\mathit{ic}}
\def\Picardfunctor{\mathrm{Pic}}
\def\Deformationcategory{\mathcal{D}\!\mathit{ef}}

% --- appended 2026-10-12 (Chief): the project own definitions of the reference/slogan/history
% display environments, copied mechanically from the shared Stacks preamble lines 8-12
% (\usepackage{verbatim} plus the three \newenvironment lines, which the chapter preamble keeps
% verbatim from the project preamble). Nothing added, removed or redefined: these environments
% carry no mathematical content and the project itself typesets them as comments. Without these
% three lines a counted statement whose own text uses them fails with Environment slogan undefined.
\usepackage{verbatim}
\newenvironment{reference}{\comment}{\endcomment}
\newenvironment{slogan}{\comment}{\endcomment}
\newenvironment{history}{\comment}{\endcomment}
\usepackage{cleveref}
\crefname{originallemma}{original lemma}{original lemmas}
\Crefname{originallemma}{Original lemma}{Original lemmas}
\makeatletter
\newlabel{constructions-lemma-proj-quasi-compact}{{01MD}{source}{}{sourcepointer-1}{}}
\newlabel{constructions-lemma-proj-quasi-compact@cref}{{[originallemma][0][]01MD}{[1][0][]source}{}{}{}}
\makeatother
\crefname{originallemma}{original lemma}{original lemmas}
\Crefname{originallemma}{Original lemma}{Original lemmas}
\makeatletter
\newlabel{constructions-lemma-relative-proj-separated}{{01O2}{source}{}{sourcepointer-2}{}}
\newlabel{constructions-lemma-relative-proj-separated@cref}{{[originallemma][0][]01O2}{[1][0][]source}{}{}{}}
\makeatother
\crefname{originallemma}{original lemma}{original lemmas}
\Crefname{originallemma}{Original lemma}{Original lemmas}
\makeatletter
\newlabel{constructions-lemma-surjective-graded-rings-map-proj}{{01N0}{source}{}{sourcepointer-3}{}}
\newlabel{constructions-lemma-surjective-graded-rings-map-proj@cref}{{[originallemma][0][]01N0}{[1][0][]source}{}{}{}}
\makeatother
\crefname{originallemma}{original lemma}{original lemmas}
\Crefname{originallemma}{Original lemma}{Original lemmas}
\makeatletter
\newlabel{constructions-lemma-surjective-graded-rings-map-relative-proj}{{07ZI}{source}{}{sourcepointer-4}{}}
\newlabel{constructions-lemma-surjective-graded-rings-map-relative-proj@cref}{{[originallemma][0][]07ZI}{[1][0][]source}{}{}{}}
\makeatother
\crefname{originallemma}{original lemma}{original lemmas}
\Crefname{originallemma}{Original lemma}{Original lemmas}
\makeatletter
\newlabel{constructions-lemma-when-invertible}{{01MS}{source}{}{sourcepointer-5}{}}
\newlabel{constructions-lemma-when-invertible@cref}{{[originallemma][0][]01MS}{[1][0][]source}{}{}{}}
\makeatother
\crefname{originallemma}{original lemma}{original lemmas}
\Crefname{originallemma}{Original lemma}{Original lemmas}
\makeatletter
\newlabel{schemes-lemma-push-forward-quasi-coherent}{{01LC}{source}{}{sourcepointer-6}{}}
\newlabel{schemes-lemma-push-forward-quasi-coherent@cref}{{[originallemma][0][]01LC}{[1][0][]source}{}{}{}}
\makeatother
\title{Every closed subscheme of a relative Proj over a quasi-compact base is itself a relative Proj of a graded quotient}
\author{The Stacks Project Authors}
\date{stacks-project:divisors.tex:local-snapshot-20261012; GNU FDL 1.2 or later}
\begin{document}
\maketitle
\noindent\textit{Editorial note.} All delivered mathematical text is the author's own, copied from the pinned Stacks Project chapter file pinned-original-source.tex; no statement or proof step was rewritten, reordered or removed.\par\medskip
\noindent\textit{Editorial note.} Following the multi-claim convention for this source, source\_url is the source-level URL https://stacks.math.columbia.edu/ and the source label lemma-closed-subscheme-proj is carried in claim\_key and here rather than in a per-claim tag URL.\par\medskip
\noindent\textit{Editorial note.} The preamble is the chapter's shared Stacks Project preamble, retained as sources/stacks-divisors/retained-preamble.tex: a line-for-line excerpt of the shared Stacks preamble staged with this snapshot. No macro was invented or redefined, and this card adds no preamble text of its own.\par\medskip
\noindent\textit{Editorial note.} Statement span 7903-7917 and proof span 7919-7983 come from the halfA pool card operations/pool-cards/stacks-divisors-halfA.json and were re-read in the pinned chapter file before this card was written.\par\medskip
\noindent\textit{Editorial note.} The proof contains the author's phrase 'Having said this' and a reference to the proof of a cited constructions lemma for the description of the ideal J'; both are preserved.\par\medskip
\noindent\textit{Editorial note.} 6 cross-references occur inside the counted spans; they are listed in external\_references as navigation pointers only.\par\medskip
\noindent\textit{Allowed standard prerequisites (editorial).}\begin{itemize}
\item The relative Proj of a quasi-coherent graded algebra is separated, and pushforward along a quasi-compact separated morphism preserves quasi-coherence, cited from the constructions and schemes chapters. Cited at proof lines 7919-7924. Makes the kernel of the canonical map of graded modules quasi-coherent, hence a legitimate graded quotient algebra.
\item A surjection of quasi-coherent graded algebras induces an everywhere defined closed immersion of relative Proj's, cited from the constructions chapter. Cited at proof lines 7926-7929. Produces the closed subscheme Z' that has to be shown equal to Z.
\item Quasi-compactness of a Proj over an affine base in terms of finitely many homogeneous elements generating the irrelevant ideal, cited from the constructions chapter. Cited at proof lines 7934-7937. Lets the proof reduce the comparison of two closed subschemes to the affine charts D\_+(f\_i).
\item Invertibility of the twists of the structure sheaf on Proj and the identification of the multiplication maps of twists of a quasi-coherent module as isomorphisms. Cited at proof lines 7951-7956. Allows homogeneous elements to be adjusted by powers of the chart equation without changing the ideal they generate.
\end{itemize}\medskip
\ifdefined\hypertarget\hypertarget{sourcepointer-1}{}\fi
\noindent\textit{Source reference.} Original lemma 01MD: focused-seeds/AGP/sources/stacks-snapshot-20261011/tags.txt, source lines 2118--2118. Same-chapter reference used in the retained author text at source line(s) 7684, 7738, 7939. It is supplied only as a navigation pointer so that the retained text compiles; the counted item delivers the author's own statement and the author's complete proof, and no part of the referenced item is copied in or substituted for the counted argument.\par\medskip
\ifdefined\hypertarget\hypertarget{sourcepointer-2}{}\fi
\noindent\textit{Source reference.} Original lemma 01O2: focused-seeds/AGP/sources/stacks-snapshot-20261011/tags.txt, source lines 2179--2179. Same-chapter reference used in the retained author text at source line(s) 7644, 7921, 8227, 8466. It is supplied only as a navigation pointer so that the retained text compiles; the counted item delivers the author's own statement and the author's complete proof, and no part of the referenced item is copied in or substituted for the counted argument.\par\medskip
\ifdefined\hypertarget\hypertarget{sourcepointer-3}{}\fi
\noindent\textit{Source reference.} Original lemma 01N0: focused-seeds/AGP/sources/stacks-snapshot-20261011/tags.txt, source lines 2141--2141. Same-chapter reference used in the retained author text at source line(s) 7960, 7990. It is supplied only as a navigation pointer so that the retained text compiles; the counted item delivers the author's own statement and the author's complete proof, and no part of the referenced item is copied in or substituted for the counted argument.\par\medskip
\ifdefined\hypertarget\hypertarget{sourcepointer-4}{}\fi
\noindent\textit{Source reference.} Original lemma 07ZI: focused-seeds/AGP/sources/stacks-snapshot-20261011/tags.txt, source lines 10148--10148. Same-chapter reference used in the retained author text at source line(s) 7931. It is supplied only as a navigation pointer so that the retained text compiles; the counted item delivers the author's own statement and the author's complete proof, and no part of the referenced item is copied in or substituted for the counted argument.\par\medskip
\ifdefined\hypertarget\hypertarget{sourcepointer-5}{}\fi
\noindent\textit{Source reference.} Original lemma 01MS: focused-seeds/AGP/sources/stacks-snapshot-20261011/tags.txt, source lines 2133--2133. Same-chapter reference used in the retained author text at source line(s) 7956. It is supplied only as a navigation pointer so that the retained text compiles; the counted item delivers the author's own statement and the author's complete proof, and no part of the referenced item is copied in or substituted for the counted argument.\par\medskip
\ifdefined\hypertarget\hypertarget{sourcepointer-6}{}\fi
\noindent\textit{Source reference.} Original lemma 01LC: focused-seeds/AGP/sources/stacks-snapshot-20261011/tags.txt, source lines 2081--2081. Same-chapter reference used in the retained author text at source line(s) 741, 6120, 7924. It is supplied only as a navigation pointer so that the retained text compiles; the counted item delivers the author's own statement and the author's complete proof, and no part of the referenced item is copied in or substituted for the counted argument.\par\medskip
\begin{lemma}
\label{lemma-closed-subscheme-proj}
Let $S$ be a scheme. Let $\mathcal{A}$ be a quasi-coherent graded
$\mathcal{O}_S$-algebra. Let
$p : X = \underline{\text{Proj}}_S(\mathcal{A}) \to S$ be the relative
Proj of $\mathcal{A}$. Let $i : Z \to X$ be a closed subscheme. Denote
$\mathcal{I} \subset \mathcal{A}$ the kernel of the canonical map
$$
\mathcal{A}
\longrightarrow
\bigoplus\nolimits_{d \geq 0} p_*\left((i_*\mathcal{O}_Z)(d)\right).
$$
If $p$ is quasi-compact, then there is an isomorphism
$Z = \underline{\text{Proj}}_S(\mathcal{A}/\mathcal{I})$.
\end{lemma}

\begin{proof}
The morphism $p$ is separated by
Constructions, Lemma \ref{constructions-lemma-relative-proj-separated}.
As $p$ is quasi-compact, $p_*$ transforms quasi-coherent modules
into quasi-coherent modules, see
Schemes, Lemma \ref{schemes-lemma-push-forward-quasi-coherent}.
Hence $\mathcal{I}$ is a quasi-coherent $\mathcal{O}_S$-module.
In particular, $\mathcal{B} = \mathcal{A}/\mathcal{I}$ is a
quasi-coherent graded $\mathcal{O}_S$-algebra. The functoriality
morphism $Z' = \underline{\text{Proj}}_S(\mathcal{B}) \to
\underline{\text{Proj}}_S(\mathcal{A})$ is everywhere defined and
a closed immersion, see Constructions, Lemma
\ref{constructions-lemma-surjective-graded-rings-map-relative-proj}.
Hence it suffices to prove $Z = Z'$ as closed subschemes of $X$.

\medskip\noindent
Having said this, the question is local on the base and we may assume
that $S = \Spec(R)$ and that $X = \text{Proj}(A)$ for some
graded $R$-algebra $A$. Assume $\mathcal{I} = \widetilde{I}$
for $I \subset A$ a graded ideal. By
Constructions, Lemma \ref{constructions-lemma-proj-quasi-compact}
there exist $f_0, \ldots, f_n \in A_{+}$ such that
$A_{+} \subset \sqrt{(f_0, \ldots, f_n)}$ in other words
$X = \bigcup D_{+}(f_i)$. Therefore, it suffices to check that
$Z \cap D_{+}(f_i) = Z' \cap D_{+}(f_i)$ for each $i$.
By renumbering we may assume $i = 0$.
Say $Z \cap D_{+}(f_0)$, resp.\ $Z' \cap D_{+}(f_0)$
is cut out by the ideal $J$, resp.\ $J'$ of $A_{(f_0)}$.

\medskip\noindent
The inclusion $J' \subset J$.
Let $d$ be the least common multiple of $\deg(f_0), \ldots, \deg(f_n)$.
Note that each of the twists $\mathcal{O}_X(nd)$ is invertible, trivialized
by $f_i^{nd/\deg(f_i)}$ over $D_{+}(f_i)$, and that for any quasi-coherent
module $\mathcal{F}$ on $X$ the multiplication maps
$\mathcal{O}_X(nd) \otimes_{\mathcal{O}_X} \mathcal{F}(m)
\to \mathcal{F}(nd + m)$ are isomorphisms, see
Constructions, Lemma \ref{constructions-lemma-when-invertible}.
Observe that $J'$ is the ideal generated by the elements $g/f_0^e$ where
$g \in I$ is homogeneous of degree $e\deg(f_0)$ (see proof of
Constructions, Lemma
\ref{constructions-lemma-surjective-graded-rings-map-proj}).
Of course, by replacing $g$ by $f_0^lg$ for suitable $l$
we may always assume that $d | e$. Then, since $g$ vanishes as a section of
$\mathcal{O}_X(e\deg(f_0))$ restricted to $Z$ we see that
$g/f_0^d$ is an element of $J$. Thus $J' \subset J$.

\medskip\noindent
Conversely, suppose that $g/f_0^e \in J$. Again we may assume $d | e$.
Pick $i \in \{1, \ldots, n\}$. Then $Z \cap D_{+}(f_i)$ is
cut out by some ideal $J_i \subset A_{(f_i)}$. Moreover,
$$
J \cdot A_{(f_0f_i)} = J_i \cdot A_{(f_0f_i)}.
$$
The right hand side is the localization of $J_i$ with respect to
$f_0^{\deg(f_i)}/f_i^{\deg(f_0)}$. It follows that
$$
f_0^{e_i}g/f_i^{(e_i + e)\deg(f_0)/\deg(f_i)} \in J_i
$$
for some $e_i \gg 0$ sufficiently divisible. This proves that
$f_0^{\max(e_i)}g$ is an element of $I$, because its restriction to each
affine open $D_{+}(f_i)$ vanishes on the closed subscheme
$Z \cap D_{+}(f_i)$. Hence $g/f_0^e \in J'$ and we conclude $J \subset J'$
as desired.
\end{proof}

\end{document}
